% STATUS: draft | reviewed | final
% PROMOTE: none | exemplars | quant-prep
% TAGS: arrays, hashing, dp, graphs, etc.
% DIFFICULTY: easy | medium | hard
% SOURCE: LeetCode XXX | Custom | Other

\ProblemTitle{Top k Frequent Elements}
\label{prob:lc-347-top-k-frequent-elements}

\ProblemSummary
Given an integer array \texttt{nums} and an integer \texttt{k}, return the \texttt{k} most frequent elements. The answer can be returned in any order.

% -----------------------------------------------------------------------------
% Optional: Author Thinking (excluded from exemplar builds)
% -----------------------------------------------------------------------------
\begin{Thinking}
Initially, my approach involved counting the frequency of each element using a hashmap (or dictionary). For every number in the array, if it hasn’t been seen before, I would add it to the hashmap with a frequency of 1. Otherwise, I would increment its frequency count.

Once the frequency count was complete, my next thought was to sort the hashmap by frequency and pick the top \texttt{k} most frequent elements. However, sorting the entire hashmap would take time proportional to the number of distinct elements, which might not be efficient enough for larger datasets.

As an optimization, I realized that I could use a \textbf{min-heap} (priority queue) to maintain the top \texttt{k} most frequent elements efficiently. This would allow me to traverse the array just once, maintaining the heap dynamically as I go, rather than sorting the entire hashmap.
\end{Thinking}

\EdgeCases
Some edge cases I considered:
\begin{itemize}
    \item If \texttt{k} is greater than the number of distinct elements in the input array.
    \item If \texttt{k} is larger than the size of the array (though it's assumed that \texttt{k} will always be a valid input).
\end{itemize}


\KeyIdea
The problem can be broken down into two main ideas:
\begin{enumerate}
    \item \textbf{Counting frequency}: We need to efficiently count how often each element appears in the input array.
    \item \textbf{Extracting the k most frequent elements}: Once we have the frequency count, we need to efficiently select the \texttt{k} most frequent elements. This can be achieved using a \textbf{min-heap} to maintain the top \texttt{k} elements as we traverse the array.
\end{enumerate}

The \textbf{hash map} helps store frequency counts, and the \textbf{min-heap} allows for efficiently keeping track of the \texttt{k} most frequent elements with the least amount of computation.

% -----------------------------------------------------------------------------
% Algorithm
% -----------------------------------------------------------------------------
\Algorithm
\begin{CodeBlock}[style=pseudo,caption={Pseudocode for Top K Frequent Elements}]{12}
1. Count the frequency of each element in \texttt{nums}.
2. Create a min-heap to store the \texttt{k} most frequent elements.
3. For each element in the frequency map:
    4. Insert the element into the heap.
    5. If the heap size exceeds \texttt{k}, remove the element with the smallest frequency.
6. Return the elements in the heap as the result.
\end{CodeBlock}

% -----------------------------------------------------------------------------
% Correctness
% -----------------------------------------------------------------------------
\Correctness
The algorithm is correct because:
\begin{itemize}
    \item We ensure every element is counted exactly once in the hashmap.
    \item By maintaining a heap of size \texttt{k}, we guarantee that after processing all elements, the heap will contain only the top \texttt{k} most frequent elements.
    \item The heap ensures that inserting and removing elements is done in logarithmic time, leading to efficient handling of the problem even with larger inputs.
\end{itemize}

% -----------------------------------------------------------------------------
% Complexity
% -----------------------------------------------------------------------------
\Complexity
\textbf{Time Complexity}: 
\begin{itemize}
    \item Counting frequencies takes \texttt{O(n)} time, where \texttt{n} is the size of the input array.
    \item Each insertion into the heap takes \texttt{O(log k)}, and since there are \texttt{n} elements, the total heap operations take \texttt{O(n log k)}.
    \item Therefore, the overall time complexity is \texttt{O(n log k)}.
\end{itemize}

\textbf{Space Complexity}: 
\begin{itemize}
    \item We store the frequencies of up to \texttt{n} elements in the hashmap, so the space complexity is \texttt{O(n)}.
\end{itemize}

% -----------------------------------------------------------------------------
% Implementation
% -----------------------------------------------------------------------------
\bigskip
\noindent\textbf{C++ Implementation}

\lstinputlisting[
  style=cpp,
  caption={C++ Solution: Top k Frequent Elements}
]{../problems/01-arrays-sequences/medium/lc-347-top-k-frequent-elements/solution.cpp}

\bigskip
\noindent\textbf{Python Implementation}

\lstinputlisting[
  language=Python,
  backgroundcolor=\color{codebg},
  basicstyle=\ttfamily\small,
  frame=single,
  breaklines=true,
  caption={Python Solution: Top k Frequent Elements}
]{../problems/01-arrays-sequences/medium/lc-347-top-k-frequent-elements/solution.py}

\ProblemDivider
